% =========================================================================
% Developed and maintained by Andrea Valsesia. 
% This project is not affiliated with "Università del Piemonte Orientale" in any way. 
% The end user is solely responsible for any unauthorized or illegal use of official university logos distributed with this template.
%
% Distributed under the MIT License.
% This project can be freely modified and redistributed as long as this copyright notice is always included in all source files.
% =========================================================================

% --- DEBUG --- %
% \usepackage{showframe}
\usepackage{lua-ul}
\usepackage{luacolor}
\usepackage{fvextra}
\fvinlineset{
    breaklines=true,
    breakanywhere=true
}
\makeatletter
\g@addto@macro\FancyVerbFormatInline{\highLight[lightgray!40]}
\makeatother

% --- LAYOUT AND DIMENSIONS ---
\usepackage[
    inner=2.5cm,
    outer=3cm,
    top=2.5cm,
    bottom=2.5cm,
    headheight=0pt,
    headsep=0pt,
    bindingoffset=0.5cm,
    footskip=1.25cm
]{geometry}
\usepackage{parskip}

% --- LANGUAGE AND TYPOGRAPHY ---
\usepackage{babel}                       % Titles of elements (e.g. "Figure 1:", "Chapter 1:", "Table 3:") are automatically translated according to target language
\usepackage{microtype}
\usepackage{lettrine}                    % The big letter at the top of chapters
\usepackage{etoolbox}
\usepackage{siunitx}                     % Professional International System units
\usepackage[version=4]{mhchem}           % Professional inline chemical formulas

\DeclareSIUnit{\rpm}{rpm}                % Setup a command for unit "revolutions per minute". Add other lines like that to add custom units

% --- FONTS ---
\usepackage{anyfontsize}
\usepackage{fontspec}                    % Use custom fonts. WARNING: The fonts used must be installed on your computer to compile the document. If the font is not found, latex will fallback to default fonts
\setmainfont{Times New Roman}
\setsansfont{Fira Sans}

% --- OTHER PACKAGES ---
\usepackage{lipsum}                      % Generates lorem ipsum paragraphs as a placeholder
\usepackage{graphicx}                    % Include images and other assets in your document
\usepackage[labelfont=bf]{caption}       % Add captions under images and plots
\usepackage{float}
\usepackage[svgnames]{xcolor}
\usepackage{amssymb}
\usepackage{comment}
\usepackage{multicol}                    % Divide the page in vertical columns (like in scientific papers)
\usepackage{tabularx}                    % Create better data tables

\usepackage{pageslts}
\usepackage{ifoddpage}                   % Check if page is odd o even in order to procedurally add different elements
\usepackage{tikz}                        % Graphic library for complex plots and virtually any image (the syntax is similar to R language)
\usetikzlibrary{calc}
\usepackage{afterpage}
\usepackage{eso-pic}

\usepackage{datetime}
\newdateformat{monthyeardate}{\monthname[\THEMONTH]\ \THEYEAR}

% --- BIBLIOGRAPHY ---
\usepackage{csquotes}
\usepackage[style=numeric, sorting=none, giveninits=true, uniquename=init, maxbibnames=99]{biblatex}
\setlength{\bibitemsep}{\baselineskip}

\usepackage{hyperref}
\pdfvariable omitcidset=1

% --- NUMBERING FOR CHAPTERS (1.) ---

\DeclareFieldFormat{labelnumberwidth}{\bfseries{#1.}}

% --- BIBLIOGRAPHY STYLE ---

\DeclareNameAlias{author}{family-given} % Format author names as: Surname, Name. (Rossi, M.)
\renewcommand{\finalnamedelim}{\addcomma\space and\space} % Add a comma before the last "and"
\renewcommand{\multinamedelim}{\addcomma\space} % Separate authors with a comma. (Rossi, M., Verdi, G.)

\DeclareFieldFormat{journaltitle}{#1}
\DeclareFieldFormat{pages}{#1}

\DeclareBibliographyDriver{article}{%
    \usebibmacro{begentry}%
    \usebibmacro{author}\addspace%
    \printtext[parens]{\printfield{year}}\addperiod\addspace%
    \printfield[titlecase]{title}\addperiod\space%
    \printfield[plain]{journaltitle}\addspace%
    \textit{\printfield{volume}}\addcomma\space%
    \setunit{}\printfield{pages}\addperiod%
    \usebibmacro{finentry}%
}

\renewbibmacro{in:}{}%
\renewcommand{\bibsetup}{\raggedright}%

\addbibresource{bibliography.bib}

% --- HEADER AND FOOTER of the page (fancyhdr) ---
\usepackage{fancyhdr}
\pagestyle{fancy}
\fancyhf{}
\fancyfoot[C]{\bfseries\sffamily\thepage}
\renewcommand{\headrulewidth}{0pt}
\renewcommand{\footrulewidth}{0pt}

\fancypagestyle{plain}{
    \fancyhead{}
}

\linespread{1.25}

% --- CHAPTER STYLES ---
\usepackage{titlesec}

\titleformat{\chapter}[display]
{}
{\fontsize{22}{0}\sffamily\selectfont\color{gray!50}\textls*[110]{CAPITOLO \ \thechapter}}  % Change this text to modify the chapter label (output: CAPITOLO 1)
{-10pt}
{\Huge\bfseries\color{black}\textsc}

\titlespacing*{\chapter}{0pc}{8pc}{3pc}

% --- SECTION STYLES ---
\titleformat{\section}
{\sffamily\selectfont\large\bfseries}
{\thesection}
{1em}
{}

\titleformat{\subsection}
{\sffamily\selectfont\normalsize\bfseries}
{\thesubsection}
{1em}
{}

\titlespacing*{\subsection}{0pt}{5ex}{2ex}

% --- CHAPTER THUMBS ---
\newbool{drawBackground}
\boolfalse{drawBackground}

\newcommand{\startThumbs}{%
    \afterpage{\global\booltrue{drawBackground}}%
}

\newcommand{\stopThumbs}{%
    \afterpage{\global\boolfalse{drawBackground}}%
}

\AddToHook{shipout/background}{%
    \ifbool{drawBackground}{%
        \ifodd\value{page}%
            \begin{tikzpicture}[remember picture, overlay]
                \coordinate (start) at ([xshift=2pt, yshift=-2.5cm-1.5cm*(\thechapter-1)]current page.north east);
                \coordinate (end) at ([xshift=-2cm, yshift=-1.5cm]start);
                \fill[black] (start) rectangle (end);
                \node[text=white, font=\huge\bfseries\sffamily] at ($(start)!0.5!(end)$) {\thechapter};
                \node[text=black, font=\small, rotate=90, anchor=east] at ($(start)!0.5!(end)$) [xshift=-1.5cm] {\rightmark};
            \end{tikzpicture}%
        \else%
            \begin{tikzpicture}[remember picture, overlay]
                \coordinate (start) at ([xshift=-2pt, yshift=-2.5cm-1.5cm*(\thechapter-1)]current page.north west);
                \coordinate (end) at ([xshift=2cm, yshift=-1.5cm]start);
                \fill[black] (start) rectangle (end);
                \node[text=white, font=\huge\bfseries\sffamily] at ($(start)!0.5!(end)$) {\thechapter};
                \node[text=black, font=\small, rotate=90, anchor=east] at ($(start)!0.5!(end)$) [xshift=-1.5cm] {\leftmark};
            \end{tikzpicture}%
        \fi%
    }{\relax}%
}
% --- END OF CHAPTER THUMBS SETUP ---

% --- SIDE MARKS ---
\renewcommand{\chaptermark}[1]{\markboth{\MakeUppercase{#1}}{}}
\renewcommand{\sectionmark}[1]{\markright{\MakeUppercase{#1}}}

% --- LETTRINE ---
\setcounter{DefaultLines}{3}
%\renewcommand{\DefaultLoversize}{1}
%\renewcommand{\DefaultLhang}{0.3}